\documentclass[11pt]{article}
\usepackage{preamble}
\setlength{\parskip}{4pt}

\begin{document}

\daybanner{AP Statistics \textperiodcentered\ Unit 2 \textperiodcentered\ Topic 2.1 \textperiodcentered\ B: 2026-10-13 \textperiodcentered\ E: 2026-10-26 \textperiodcentered\ TEACHER KEY}{Counts or Conditional Percents?}

\sectionbanner{CED TETHER}

{\small
\begin{itemize}[leftmargin=1.5em,itemsep=2pt,topsep=2pt]
  \item \textbf{VAR-1.D} Identify questions to be answered about possible relationships in data. \textbf{[Skill 1.A]}
  \item \textbf{VAR-1.D.1} Apparent patterns and associations in data may be random or not.
\end{itemize}
}
\textbf{Essential question:} When group sizes differ, which comparison reveals whether two categorical variables are related?\par
\textbf{DOK-3 skill:} 4.B \textperiodcentered\ \textbf{FRQ pattern:} {\small\texttt{counts-\allowbreak{}vs-\allowbreak{}conditional-\allowbreak{}percents-\allowbreak{}are-\allowbreak{}the-\allowbreak{}variables-\allowbreak{}related}}\par
\textbf{DOK rationale:} Part (c) coordinates unequal group sizes, two conditional percentages, and the limits of observational data to adjudicate competing claims about relationship and cause.\par

\frameworkphaseheader{Do Now (first take) $\rightarrow$ Explore (video + follow-along) $\rightarrow$ Exit (finish + turn in)}{1 $\rightarrow$ 3}{45}{%
\begin{itemize}[leftmargin=1.5em,itemsep=2pt,topsep=2pt]
  \item Board slide up at the bell. Read Kai's and Rosa's statements without confirming either.
  \item Begin the video follow-along at minute 5.
  \item Return to the table at minute 33. Circulate and ask students to name every denominator they use.
  \item Collect at the door. Score part (c) only, E/P/I.
\end{itemize}
}{%
\begin{itemize}[leftmargin=1.5em,itemsep=2pt,topsep=2pt]
  \item Choose the more useful line of reasoning and cite one table detail.
  \item Complete the video follow-along about variable types, displays, and possible relationships.
  \item Finish (a)--(c), revise the first take in the exit line, and turn in the sheet.
\end{itemize}
}{%
\begin{itemize}[leftmargin=1.5em,itemsep=2pt,topsep=2pt]
  \item What question does the count 75 answer? What question would 75 divided by 150 answer?
  \item If both sport groups had 100 students, what counts would these percentages produce?
  \item Does knowing sport participation change your prediction of the sleep category?
  \item What assignment or control would be needed before using the word causes?
\end{itemize}
}{Solo teacher. Circulate ~5 pairs. Require both conditional percentages and the unequal-row-total mechanism before accepting a relationship claim; do not announce which student is right.}

\begin{callout}{\IconWarn\ WATCH FOR}{calloutred}
\begin{itemize}[leftmargin=1.5em,itemsep=2pt,topsep=2pt]
  \item Dividing each cell by 200 instead of conditioning within the sport-participation row.
  \item Treating a larger raw count as a larger within-group chance.
  \item Changing associated into caused without evidence from an experiment.
\end{itemize}
\end{callout}

\sectionbanner{SCORING PART (c) --- E / P / I}

\textbf{Expected elements}
\begin{itemize}[leftmargin=1.5em,itemsep=2pt,topsep=2pt]
  \item \textbf{decides-from-conditional-distributions} (required) --- States that the variables appear related and bases the decision on within-group percentages
  \item \textbf{cites-both-percentages} (required) --- Cites 70 percent for athletes and 50 percent for nonathletes
  \item \textbf{explains-count-misreading} (required) --- Explains that 75 versus 35 reflects row totals of 150 versus 50 and can reverse the likelihood comparison
  \item \textbf{bounds-causal-claim} (required) --- States that observational association does not establish that sport participation caused a sleep difference
\end{itemize}

{\small\renewcommand{\arraystretch}{1.25}
\noindent\begin{tabular}{|p{0.5in}|p{5.9in}|}
\hline
\textbf{E} & Decides association from both conditional percentages, explains how unequal group sizes make the raw-count comparison misleading, and correctly rejects a causal conclusion. \\ \hline
\textbf{P} & Reaches the supported association decision and cites at least one relevant percentage but omits the unequal-size mechanism, the second percentage, or the causal limit. \\ \hline
\textbf{I} & Uses 75 versus 35 as the likelihood comparison, says the variables are unrelated despite the conditional difference, or claims sport participation caused the sleep pattern. \\ \hline
\end{tabular}}

\textbf{Common mistakes}
\begin{itemize}[leftmargin=1.5em,itemsep=2pt,topsep=2pt]
  \item Dividing 35 and 75 by the grand total instead of their row totals
  \item Saying 75 is a larger percentage because it is a larger count
  \item Claiming that association proves playing a sport increases sleep
\end{itemize}

\clearpage
\focusbanner{TODAY'S PROBLEM --- annotated}

Suppose a counselor records whether each of 200 students plays a school sport and whether the student usually gets at least 8 hours of sleep on a school night.\par\smallskip \textbf{Kai:} \emph{Seventy-five nonathletes but only 35 athletes get at least 8 hours, so playing a sport goes with less sleep.}\par \textbf{Rosa:} \emph{Those counts compare groups of different sizes; 35 out of 50 may tell a different story from 75 out of 150.}\par
\begin{center}
\textbf{Usual sleep on a school night}\par\smallskip
\begin{tabular}{|l|c|c|c|}
\hline
\textbf{Sport?} & \textbf{$\geq 8$ h} & \textbf{$<8$ h} & \textbf{Total} \\ \hline
\textbf{Plays} & 35 & 15 & 50 \\ \hline
\textbf{Does not} & 75 & 75 & 150 \\ \hline
\textbf{Total} & 110 & 90 & 200 \\ \hline
\end{tabular}
\end{center}
\begin{firsttakebox}{FIRST TAKE (5 min)}
Which student's reasoning seems more useful for deciding whether sport participation and sleep are related? Name one table detail.\par
\answer{Not graded. Expect the 75-to-35 count comparison to attract students before they normalize for group size.}
\end{firsttakebox}

\textit{Rules callout ``COMPARE WITHIN GROUPS (keep on the page)'' is printed on the student sheet.}\par\medskip

\dokbadge{1}~\textbf{(a)} Identify the two categorical variables and list the categories of each.\par
\answer{Variables: sport participation (plays a school sport, does not play) and usual school-night sleep category (at least 8 hours, under 8 hours).}

\dokbadge{2}~\textbf{(b)} Compute the conditional distribution of sleep category within each sport-participation group. Compare the percentages who get at least 8 hours.\par
\answer{Among athletes, $35/50=70\%$ get at least 8 hours and $15/50=30\%$ get under 8 hours. Among nonathletes, $75/150=50\%$ get at least 8 hours and $75/150=50\%$ get under 8 hours. The at-least-8-hours percentages differ by 20 percentage points.}

\dokbadge{3}~\textbf{(c)} Decide whether sport participation and sleep category appear related for these students. Cite the two conditional percentages that settle your decision. Explain what Kai's comparison of 75 with 35 would mislead a reader into believing and why the unequal row totals create that impression. Finally, state what this table cannot establish about playing a sport causing a change in sleep.\par
\answer{The variables appear related for these 200 students because the conditional sleep distributions differ: $70\%$ of athletes versus $50\%$ of nonathletes usually get at least 8 hours. Kai's 75-versus-35 comparison could make a reader think nonathletes are more likely to get enough sleep, but there are three times as many nonathletes overall (150 versus 50), so raw counts answer how many, not which group has the larger within-group rate. The table shows association only. Because students were observed rather than assigned to play a sport, other differences between the groups could explain the pattern; it cannot establish that sport participation caused more sleep.}

\begin{summaryexitbox}{EXIT}
One thing the video changed about my first take: \hrulefill
\end{summaryexitbox}

\end{document}
